\section{Experiencia}
\outerlist{

% \entrybig
% 	{\textbf{Ayesa}}{Bilbao}
% 	{Científico de datos en el Instituto de Ibermática de Innovación (i3B).}{Feb 2024 - Ahora}
    
%     \text{Participación en proyectos de I+D de inteligencia artificial aplicada de escala nacional y europea. Tareas de investigación en nuevos algoritmos y propuestas innovadoras, análisis y procesamiento de datos, experimentación, modelado predictivo, despliegue y comunicación de resultados.}
    
% \innerlist{
% 	\entry{Desarrollo de modelos de machine learning para predecir la concentración de níquel en procesos industriales de galvanizado. Aplicación de técnicas de IA explicable (XAI) para interpretar los resultados y facilitar la toma de decisiones. Comunicación de hallazgos a perfiles técnicos y no técnicos del sector industrial.}
%         \entry{Desarrollo de un marco basado en Machine Learning para la detección temprana de desviaciones, retrasos y cuellos de botella en proyectos de construcción. Aplicación de técnicas de clustering, selección de variables, detección de anomalías y XAI para optimizar la toma de decisiones en la industria AEC (Arquitectura, Ingeniería y Construcción).}
%         \entry{Modelado predictivo y process mining aplicado a producción de acero: clasificación binaria, análisis de variables, métricas de conformidad (fitness, precision) y visualización de flujos con PM4Py. Implementación de pipelines MLOps con MLflow para seguimiento de experimentos y fine-tuning de modelos. Generación de datos sintéticos con VAEs y GANs.}
%         \entry{Predicción de series temporales de consumo energético, incluyendo análisis y limpieza de datos, detección de outliers, modelado con técnicas clásicas de ML (ARIMA, Prophet, LSTM) y evaluación comparativa de modelos de predicción a corto y medio plazo.}
% 	% \entryextra{A more detailed explanation of the project}
% }

\entrybig
	{\textbf{Instituto Ibermática de Innovación (Ayesa)}}{Bilbao}
	{Científico de datos}{Feb 2024 - Ahora}
    
    \par{\small Participación en proyectos de I+D en IA aplicada a la industria a nivel nacional y europeo, llevando a cabo investigaciones sobre nuevos algoritmos, análisis de datos, experimentación, creación de modelos predictivos, visualizaciones y comunicación de resultados.}
    
% \innerlist{
% 	\entry{Desarrollo de modelos de machine learning para la predicción de la concentración de níquel en procesos industriales de galvanizado. Aplicación de técnicas de IA explicable (XAI) para interpretar los resultados y facilitar la toma de decisiones. Comunicación de hallazgos a perfiles técnicos y no técnicos del sector industrial.}
        
%     \entry{Diseño de un marco basado en Machine Learning (ML) para la detección temprana de desviaciones, retrasos y cuellos de botella en proyectos de construcción. Aplicación de clustering, selección de variables, detección de anomalías y XAI para optimizar la toma de decisiones en la industria AEC (Arquitectura, Ingeniería y Construcción). Integración de modelos en entornos productivos.}
        
%     \entry{Modelado predictivo y process mining aplicado a producción de acero: clasificación binaria, análisis de variables, métricas de conformidad (fitness, precision) y visualización de flujos con PM4Py. Implementación de pipelines MLOps con MLflow para el seguimiento de experimentos, automatización del ciclo de vida de modelos y fine tuning de hiperparámetros. Generación de datos sintéticos con modelos generativos como VAEs y GANs para mejorar la robustez de los modelos.}
        
%     \entry{Predicción de series temporales de consumo energético, incluyendo análisis y limpieza de datos, detección de outliers, modelado con técnicas como ARIMA, Prophet, modelos de boosting y LSTM, y evaluación comparativa de modelos. Participación en el diseño y validación de componentes software para la integración de modelos predictivos en sistemas industriales.}
% }

\innerlist{
    % \entry{Desarrollo de modelos de ML para la predicción de la concentración de níquel en procesos industriales de galvanizado, alcanzando F1-score de 0.99 mediante el uso de variables industriales. Aplicación de técnicas de IA explicable (XAI) para identificar drivers clave y mejorar la interpretabilidad del modelo en la toma de decisiones.}

    \entry{Diseño de un framework de ML para la detección temprana de riesgos en documentos asociados a proyectos de construcción, dando lugar a una publicación científica. Obtención de un F1-score de 0.73 en la detección de cuellos de botella mediante reglas interpretables, identificación de un 6\% de anomalías y segmentación en 3 clusters de comportamiento de proyectos. Aplicación de clustering, modelos ensemble de detección de anomalías y XAI para la identificación de variables críticas y patrones de riesgo en la industria AEC.}
        
    \entry{Desarrollo de una herramienta en entorno cloud para visualización interactiva de rutas teóricas vs reales en producción de acero, incluyendo análisis y trazabilidad de producto. Implementación de dashboards basados en PM4Py para el estudio de coincidencias y desviaciones de rutas. Implementación de pipelines MLOps con MLFlow para seguimiento de experimentos con fine-tuning de modelos en tareas de clasificación binaria. Generación de datos sintéticos con VAEs y GANs para datos mixtos, mejorando la robustez de los modelos.}
        
    \entry{Desarrollo de un sistema end-to-end de forecasting horario de consumo eléctrico, desde ingestión de datos hasta despliegue vía API REST. Comparativa de modelos de regresión con optimización de hiperparámetros y validación temporal. Mejor modelo obtiene reducción de ~77\% frente a baseline naive. Interpretabilidad con SHAP y despliegue con FastAPI y MLflow.}

    \entry{Investigación comparativa de ML clásico vs. ML cuántico (QML) en datasets ad-hoc y datasets de clasificación y regresión molecular, implementando pipelines con SKLearn, Qiskit y MLflow. Entrenamiento, evaluación, análisis de robustez y escalabilidad para modelos clásicos frente a sus análogos cuánticos.}

    \entry{Desarrollo de un sistema de predicción multivariable en tiempo real para una máquina CNC industrial, procesando datos de sensores a 1 Hz y generando predicciones de 1 a 5 pasos para 7 variables de proceso. Implementación de un ensemble de modelos basado en online learning y batch learning con ponderación dinámica por MAE inverso. Diseño de un registro de modelos por herramienta para aislar el estado de los modelos ante los cambios de herramienta, mejorando así la adaptación a señales industriales no estacionarias.}
}



\entrybig
	{\textbf{Hubbell}}{Bilbao}
	{Científico de datos}{Oct 2023 - Feb 2024}
    
% \innerlist{
	\par{\small Desarrollo de mi Trabajo de Fin de Máster en Ciencia de Datos: "Advanced Metering Infrastructure-oriented data imputation through machine learning techniques", con el objetivo de estudiar e imputar valores faltantes en series temporales de smart meters mediante modelos de ML aplicados a datos reales del sector eléctrico.}
	% \entryextra{A more detailed explanation of the project}
% }


\entrybig
	{\textbf{EDP Energy}}{Bilbao}
	{Científico de datos}{Abr 2022 - Jul 2022}
    
% \innerlist{
	\par{\small Prácticas extracurriculares en el departamento de Inteligencia de negocio y Big Data de EDP Energía. 
	% Creación y optimización de modelos de ML aplicado a negocios comerciales de energía e investigación de nuevos sistemas Open Data.
    % Durante las prácticas me dediqué a la investigación y procesamiento de sistemas Open Data y trabajé en proyectos de machine learning aplicados a negocios comerciales de energía. Las tecnologías que utilicé fueron Python, SAS y PowerBI.
 }
	% \entryextra{A more detailed explanation of the project}
% }

}

\clearpage